% main.tex
% Modular Document Example - Demonstrates \include for large projects
% KU Leuven - Faculteit Industriële Ingenieurswetenschappen
%
% This template shows how to split a large document into multiple files.
% Use this structure for thesis, large summaries, or multi-chapter documents.

\documentclass[a4paper,11pt]{report}
\usepackage{../../school-macros}

% ============================================================================
%                           DOCUMENT METADATA
% ============================================================================
% ============================================================================
%                           DOCUMENT METADATA
% ============================================================================
\newcommand{\titel}{Ontwerp van een industriële sturing}
\newcommand{\auteur}{Ruben Ryckaert, Andere Auteur}
\newcommand{\academiejaar}{Academiejaar 2025-2026}

% Optional: Define course info for headers
\newcommand{\vaknaam}{Voorbeeldvak}

% ============================================================================
%                              DOCUMENT
% ============================================================================
\begin{document}

% --- Title Page (KU Leuven style) ---
\kuleuventitle
    {\titel}
    {\vaknaam}
    {\auteur}
    {\academiejaar}

 
% --- Front Matter ---
    \begin{abstract}
Deze samenvattingen zijn beschikbaar op GitHub waar je kunt bijdragen aan verbeteringen: \url{https://github.com/KUL-Industriele-ingenieurs/Samenvattingen-Ku-leuven-Industriele-ingenieurs}

Je bent helemaal vrij en ik raad je zelf aan om info toe te voegen aan de documenten als je de uitleg niet goed vindt.
Je leert je vak en je raakt dan ook bekend met LaTeX wat handig is voor later.
\end{abstract}
\tableofcontents

% ============================================================================
%                           CHAPTER INCLUDES
% ============================================================================
% Use \include{} for chapters - each starts on a new page
% Use \input{} for smaller sections that shouldn't force page breaks
%
% Benefits of \include:
% - LaTeX can skip unchanged chapters (faster compilation)
% - Each chapter can be compiled separately for testing
% - Better organization for large projects
%
% Tip: Use \includeonly{chapter1,chapter2} in preamble to compile 
%      only specific chapters while keeping all cross-references

% Symbol list and formularium are auto-generated

\include{labo-info.tex}
\chapter{Introductie}
\label{chap:introductie}

\begin{wrapfigure}{r}{0.5\textwidth}[10]
  \centering
  \includegraphics[width=0.6\textwidth]{image1.png}
  \caption{}
  \label{fig:image1.png}
 \end{wrapfigure}

We leggen deze figuur uit met een eenvoudig voorbeeld.
Je hebt een doodgewone knop die je kan indrukken.
Jij als \textbf{operator} geef controle inputs aan de knop, 
bijvoorbeeld door hem in te drukken.

\textbf{De logic controller (PLC) }ontvangt deze input en 
verwerkt deze volgens een vooraf gedefinieerd programma.

De knop krijgt info door sensors die bijvoorbeeld info geven of 
de knop al dan niet is ingedrukt van de machine

Als de knop is ingedrukt gaat door er een \textbf{logic actions} gebeuren waarbij de \textbf{machine} een actie uitvoert,
bijvoorbeeld een motor starten of een lamp laten branden.
en de \textbf{operator} krijgt feedback van de \textbf{machine}, bijvoorbeeld door een lamp die gaat branden of een geluidssignaal.


% ============================================================================
%                              APPENDICES
% ============================================================================
\appendix



\end{document}
